\subsection{Separating transcendence bases}

DEFINITION 1. --- Let L be an extension of a field K. A transcendence basis B of L over K is said to be separating if the algebraic extension L of K(B) is separable.

If K has characteristic 0, every transcendence basis of L over K is separating because every algebraic extension of a field of characteristic 0 is separable (V, p.~37, Cor.). If an extension admits a separating transcendence basis, it is separable (V, p.~121, Prop. 6 and p.~122, Prop. 9). The following theorem shows that every finitely generated separable extension admits a separating transcendence basis; this restriction is essential (V, p.~177, Ex. 1).

THEOREM 5. --- Let K be a field, L an extension of K and \((x_{i})_{i\in I}\) a family of elements of L. If the family \((x_{i})_{i\in I}\) is a separating transcendence basis of L over K, then the family \((dx_{i})_{i\in I}\) is a basis of the vector space \(\Omega_{\mathrm{K}}(\mathrm{L})\) over L. The converse holds if L is a finitely generated separable extension of K.

Put \(\mathbf{M} = \mathbf{K}(x_i)_{i \in I}\) and denote by \(\alpha\) the canonical mapping of \(\Omega_{\mathbf{K}}(\mathbf{M}) \otimes_{\mathbf{M}} \mathbf{L}\) into \(\Omega_{\mathbf{K}}(\mathbf{L})\) . If \((x_i)_{i \in I}\) is a separating transcendence basis of L over K, the family \((d_{\mathbf{M}/\mathbf{K}} x_i)_{i \in I}\) is a basis of the vector M-space \(\Omega_{\mathbf{K}}(\mathbf{M})\) (V, p.~128, Prop. 3) and \(\alpha\) is an isomorphism of vector L-spaces since L is algebraic and separable over M (V, p.~130, Cor. 2). Since \(\alpha(d_{\mathbf{M}/\mathbf{K}} x_i \otimes 1) = d_{\mathbf{L}/\mathbf{K}} x_i\) , the family \((d_{\mathbf{L}/\mathbf{K}} x_i)_{i \in I}\) is thus a basis of \(\Omega_{\mathbf{K}}(\mathbf{L})\) over L.

Conversely, suppose that L is a separable finitely generated extension of K and that the family \((d_{L/K}x_{i})_{i\in I}\) is a basis of the vector space \(\Omega_{K}(L)\) over L. By Cor. 1 of V, p.~134, the transcendence degree of L over K is equal to the dimension of \(\Omega_{K}(L)\) over L, hence to the cardinal of I. From the exact sequence \((\mathrm{E}_{\mathrm{K},\mathrm{M},\mathrm{L}})(\mathrm{V},\mathrm{p}.128)\) we have \(\Omega_{\mathrm{M}}(\mathrm{L})=0\) ; since L is a finitely generated extension of M, Cor. 2 of V, p.~135 shows that L is algebraic and separable over \(\mathbf{M}=\mathbf{K}(x_{i})_{i\in I}\) ; since the transcendence degree of L over K is finite and equal to the cardinal of I, the family \((x_{i})_{i\in I}\) is a transcendence basis of L over K (V, p.~110, Cor. 1).

COROLLARY. --- Let L be a separable finitely generated extension of K and let S be a subset of L such that \(L = K(S)\) . Then there exists a separating transcendence basis B of L over K, contained in S.

Since \(\Omega_{\mathrm{K}}(\mathrm{L})\) is generated by the differentials of the elements of S, there exist s elements \(x_{1}, \ldots, x_{s}\) of S such that \((dx_{1}, \ldots, dx_{s})\) is a basis of \(\Omega_{\mathrm{K}}(\mathrm{L})\) over L. Now it suffices to apply Th. 5.

Remark. --- Let L be a separable finitely generated extension of a field K of characteristic \(p \neq 0\) ; there may exist transcendence bases of L which are not separating. It is enough to observe that \(\{X^{p}\}\) is a transcendence basis of \(\mathrm{K}(\mathbf{X})\) but \(\mathrm{K}(\mathbf{X})\) is a p-radical extension of degree p of \(\mathrm{K}(\mathbf{X}^{p})\) .

PROPOSITION 5. --- Let L and M be two extensions of a field K contained in a given extension and algebraically disjoint over K. If M is separable over K, then L(M) is separable over L.

It is enough to show that for every finite subset S of M, L(S) is separable over L (V, p.~122, Prop. 8). Let S be a finite subset of M. Since the field K(S) is separable over K, it has a separating transcendence basis B (Cor. of Th. 5). Since L and M are algebraically disjoint over K, B is a transcendence basis of L(B) over L (V, p.~113, Prop. 12). Further, every element of S is algebraic and separable over K(B), hence over L(B) (V, p.~39, Cor. 2). We deduce (V, p.~39, Prop. 6) that L(S) = L(B)(S) is algebraic and separable over L(B), hence L(S), is separable over L.

COROLLARY. --- If L and M are separable extensions, algebraically disjoint over K, then the field \(\mathbf{K}(\mathbf{L} \cup \mathbf{M})\) is separable over K.

For \(K(L \cup M) = L(M)\) is separable over L by Prop. 5 (because M is separable over K) and L is separable over K, whence the corollary (V, p.~122, Prop. 9).

\pandocbounded{\includegraphics[keepaspectratio]{images/part002_2c83366d.jpg}}

The hypothesis that the extensions L and M are algebraically disjoint is indispensable in Prop. 5 and its corollary. For let K be an imperfect field of characteristic \(p \neq 0\) and E an extension of the form \(\mathrm{K}(x, a)\) with x transcendental over K and a p-radical of height 1 over K; put \(\mathrm{L} = \mathrm{K}(x)\) and \(\mathrm{M} = \mathrm{K}(x + a)\) . Then \(x + a\) is transcendental over K (if not, \(x = (x + a) - a\) would be algebraic over K) and the fields L and M are separable over K. However, \(\mathrm{K}(\mathrm{L} \cup \mathrm{M}) = \mathrm{K}(x, a)\) is p-radical of degree p over \(\mathrm{L} = \mathrm{K}(x)\) and is not separable over L, nor over K.

